\section{Main Theorem}

\begin{theorem}[Fixed-span positive limiting objective]\label{thm:main}
Under Assumptions~\ref{assump:dimension}--\ref{assump:joint_initialization},
set
\[
 r_0=1,
 \qquad
 \alpha=\frac14,
 \qquad
 L(r)=r^{1+\alpha}=r^{5/4}.
\]
For every integer $r\ge r_0$, every integer $n\ge8L(r)$, every integer
$k$ with $r<k\le L(r)$, and every deterministic base triple allowed by
Assumption~\ref{assump:arbitrary_base}, run the two constrained methods above
with their independent Gaussian starts and their shared smoothed tensor $T$.
Then, with probability at least $1/4$ over the smoothing variables and both
complete initialization triples, simultaneously for
$M\in\{\mathrm{cALS},\mathrm{cGD}\}$,
\[
 \|(I-P_{\mathcal H_M})T\|_F^2\ge\frac34\|T\|_F^2,
 \qquad
 F_M(t)\ge\frac38\|T\|_F^2\quad\text{for every }t\ge0,
\]
and the finite scalar limit satisfies
\[
 \lim_{t\to\infty}F_M(t)\ge\frac38\|T\|_F^2.
\]
Equivalently,
\[
 \mathbb P\!\left[
 \bigcap_{M\in\{\mathrm{cALS},\mathrm{cGD}\}}
 \left\{
 \lim_{t\to\infty}F_M(t)\text{ exists as a finite real number and }
 \lim_{t\to\infty}F_M(t)\ge\frac38\|T\|_F^2
 \right\}
 \right]\ge\frac14.
\]
All displayed constants are exact and hide no dependence on
$r,n,k,q,\rho,T$, the deterministic bases, either initialization, the method,
or the iteration index.  The probability statement is pointwise uniform over
the deterministic base triple, not a simultaneous event over all triples.
The conclusion is only for the two displayed fixed-one-mode constrained
procedures and does not extend to ordinary unconstrained ALS or full-variable
gradient descent.
\end{theorem}
